\chapter{Huber Joint Boundaries and the Sign Information Bridge}
\label{sec:p15-huber-joint-endpoints}

When scale is estimated jointly with location, an atom at the clipping
point can move wholly into or out of the active observations under
an arbitrarily small scale change. The limit is then governed by
directional, rather than ordinary linear, score derivatives.
This chapter supplies that limit for the equations in
\cref{sec:p15-huber-joint-proposals}, completes the normal-mixture
consistency bridge, and proves the scale minimax comparison when
estimators can use signs as well as magnitudes.
The original equations and variance discussion are
\cite[Secs.~10--11, pp.~94--100]{Huber1964RobustLocation}.
The nonsmooth limit and its consequences below are a proved
completion of that discussion.

\section{Identification before a limiting distribution}
\label{sec:p15-huber-joint-identification}

\begin{hbdefinition}{def:p15-huber-joint-boundary-model}{Symmetric laws and the joint equations}
Fix \(k>0\), \(0<\beta<k^2\), and a symmetric Borel probability
\(F\) on \(\mathbb R\). Write
\[
 \psi(u)=\max(-k,\min(u,k)),\qquad
 \rho(u)=
 \begin{cases}u^2/2,&|u|\leq k,\\ k|u|-k^2/2,&|u|>k.\end{cases}
\]
Observations are i.i.d.\ with law \(F\). The joint equations are
\[
 \frac1n\sum_i\psi((X_i-T)/S)=0,\qquad
 \frac1n\sum_i\psi((X_i-T)/S)^2=\beta,\qquad S>0.
\]
Their finite criterion is
\[
 L_n(t,s)=\frac1n\sum_i
       \{s\rho((X_i-t)/s)+\beta s/2\}.
\]
No moment of \(X\) is assumed. Population criterion differences
are defined by subtracting the same summand at a fixed positive-scale
reference point before taking expectation.
To define the asymptotic sequence on every sample, use the measurable
rational-approximation selection in the annex and fill an undefined
value with \((0,1)\). The selection is the unique positive root
eventually almost surely. No equivariance of the exceptional
extension is asserted.
\end{hbdefinition}

\begin{hbtheorem}{thm:p15-huber-joint-identified}{proved=annex:p15.huber.joint.identified}{The exact centered scale condition and eventual unique estimation}
Put \(m_0=F\{0\}\). A unique positive centered scale \(S_0\) satisfying
\[
                    E_F\psi(X/S_0)^2=\beta
\]
exists exactly when
\[
                         m_0<1-\beta/k^2.
\]
Under this strict condition the centered population criterion
has the unique joint minimizer \((0,S_0)\).
The sample equations have a unique positive-scale solution eventually
almost surely, and that solution tends almost surely to \((0,S_0)\).
Thus this conclusion does not require the largest atom to be
smaller than \(1-\beta/k^2\), nor a positive continuous component.
It is a population/asymptotic conclusion; the earlier finite-sample
multiplicity condition remains a useful sufficient condition
for each particular sample.

If \(m_0>1-\beta/k^2\), no positive centered scale solves the equation.
At equality there is either no positive solution or a nonunique
interval extending down to zero. For example, with
\(r=\beta/k^2\),
\[
 F=(1-r)\delta_0+\tfrac r2(\delta_{-1}+\delta_1)
\]
has every \(0<S\leq1/k\) as a centered solution.
Replacing the two nonzero atoms by a symmetric uniform law on
\([-1,1]\), still of total mass \(r\), gives no positive solution.
\end{hbtheorem}

\section{The complete directional equation}
\label{sec:p15-huber-joint-directional}

\begin{hbdefinition}{def:p15-huber-joint-directional-constants}{Interior information and clipping atoms}
Assume the strict identification condition. Put
\[
 U=X/S_0,\quad q=kS_0,\quad
 A=P(|U|<k),\quad C=E[U^2{\bf1}_{\{|U|<k\}}],
 \quad p=P(U=k)=P(U=-k),
\]
and
\[
 \eta^2=\tfrac14\{\operatorname{Var}(\psi(U)^2)\}.
\]
Here \(A,C>0\). An increment \(h=(a,b)\) parameterizes
\(t=S_0a\), \(s=S_0(1+b)\) when \(b>-1\).
Every \(h\in\mathbb R^2\) is an allowed direction after multiplication
by a sufficiently small positive number.
Let \(x_+=\max(x,0)\), \(v_+=(1,k)^{\mathsf T}\),
\(v_-=(-1,k)^{\mathsf T}\), and define
\[
 \begin{split}
 \Gamma(a,b)&=\tfrac12Aa^2+\tfrac12Cb^2+
       \tfrac p2\{(a+kb)_+^2+(-a+kb)_+^2\},\\
 \mathcal D(h)&=\nabla\Gamma(h).
 \end{split}
\]
Explicitly,
\[
 \begin{split}
 \mathcal D_1(a,b)&=Aa+
       p\{(a+kb)_+-(-a+kb)_+\},\\
 \mathcal D_2(a,b)&=Cb+
       kp\{(a+kb)_++(-a+kb)_+\}.
 \end{split}
\]
\end{hbdefinition}

\begin{hbtheorem}{thm:p15-huber-joint-directional}{proved=annex:p15.huber.joint.directional}{Exact population derivatives at arbitrary clipping atoms}
Define the two population moments by
\[
 m(h)=E\left[
 \begin{pmatrix}
 \psi((U-a)/(1+b))\\
 \{\psi((U-a)/(1+b))^2-\beta\}/2
 \end{pmatrix}\right],\qquad b>-1.
\]
Then \(m(0)=0\), and, as \(t\downarrow0\),
\[
                 \frac{m(th)-m(0)}t\longrightarrow-\mathcal D(h),
\]
uniformly for \(h\) in any compact set. The same statement holds
when \(h\) varies and converges to a fixed direction.
The two atoms contribute the positive-part terms displayed above.
When \(p>0\), this derivative map is not linear.
\end{hbtheorem}

\begin{hbtheorem}{thm:p15-huber-joint-cone-inverse}{proved=annex:p15.huber.joint.cone.inverse}{An everywhere defined inverse with four explicit pieces}
The map \(\mathcal D:\mathbb R^2\to\mathbb R^2\) is one-to-one
and onto. Its inverse is Lipschitz with constant
\(1/\min(A,C)\). To compute \(h=\mathcal D^{-1}(z)\), consider
the four choices \(e_+,e_-\in\{0,1\}\), form
\[
 M_e=
 \begin{pmatrix}A&0\\0&C\end{pmatrix}
       +p e_+v_+v_+^{\mathsf T}+p e_-v_-v_-^{\mathsf T},
\]
and solve \(h=M_e^{-1}z\). Retain a piece exactly when
\(e_+=1\) corresponds to \(a+kb\geq0\), and \(e_-=1\)
to \(-a+kb\geq0\), with nonpositive inequalities for zero choices.
On shared boundaries the retained values agree.
All four matrices are positive definite; this is not an algorithm
requiring an unproved inverse or an unspecified choice of root.
\end{hbtheorem}

\begin{hbtheorem}{thm:p15-huber-joint-atomic-limit}{proved=annex:p15.huber.joint.atomic.limit}{The full joint limit, including the atoms}
Under the strict identification condition, let
\[
 H_n=\sqrt n\,(T_n/S_0,\ S_n/S_0-1)^{\mathsf T}.
\]
Then
\[
 H_n\Longrightarrow H=\mathcal D^{-1}(Z),\qquad
 Z_1\sim N(0,\beta),\quad Z_2\sim N(0,\eta^2),
\]
where the two Gaussian coordinates are independent.
The theorem allows arbitrary atoms at the two clipping points
and everywhere else, subject only to the stated identification
condition. It also allows \(\eta^2=0\).
The proof supplies the root rate and the uniform local empirical
expansion, rather than assuming either one.
\end{hbtheorem}

\begin{hbtheorem}{thm:p15-huber-joint-law-and-variance}{proved=annex:p15.huber.joint.law.variance}{Gaussian special case, nonsmooth law and the proposed variance estimate}
If \(p=0\), the joint limit is Gaussian, with independent
coordinates of variances \(\beta/A^2,\eta^2/C^2\);
the second coordinate may be identically zero.
If \(p>0\), then \(\eta^2>0\) and the joint limit is not Gaussian.
On the interior of each parameter cone its density is
\[
 f_H(h)=\frac{\det M_e}{2\pi\sqrt{\beta\eta^2}}\,
      \exp\left(-\tfrac12(M_eh)^{\mathsf T}
                  \begin{pmatrix}\beta&0\\0&\eta^2\end{pmatrix}^{-1}
                         (M_eh)\right).
\]
The location limit is symmetric about zero, and
\[
 \frac{\beta}{(A+2p)^2}\leq E H_1^2\leq\frac{\beta}{A^2};
\]
both inequalities are strict when \(p>0\).
The pointwise limiting truncated risk for location is
\[
 \lim_{R\to\infty}\lim_{n\to\infty}
       E\min(nT_n^2,R^2)=S_0^2EH_1^2 .
\]
This is not a claim of convergence of untruncated finite-sample
variances.

If \(p>0\), the strict-interior version of the proposed variance
estimate in \cref{thm:p15-huber-joint-regular-limit} instead has
the random limit
\[
 n\widehat v_n\Longrightarrow
 \frac{S_0^2\beta}
 {[A+p{\bf1}_{\{H_1+kH_2>0\}}
     +p{\bf1}_{\{-H_1+kH_2>0\}}]^2}.
\]
Each of the three possible denominator values has positive
probability. It consequently does not converge in probability
to a constant variance at these boundary laws.
Set that estimate to zero on its undefined finite-sample events;
their probabilities tend to zero and do not change this limit.
\end{hbtheorem}

\begin{hbtheorem}{thm:p15-huber-joint-worked-atomic-law}{proved=annex:p15.huber.joint.worked.atomic.law}{An exactly calculated four-point law}
Take \(k=1\), \(\beta=5/8\), and give each of
\(-1,-1/2,1/2,1\) probability \(1/4\). Then \(S_0=1\),
\[
 A=1/2,\quad C=1/8,\quad p=1/4,\quad \eta^2=9/256.
\]
For \(r=|z_1|\) and \(w=z_2\), the absolute location coordinate
of \(\mathcal D^{-1}(z)\) is
\[
 |H_1|=
 \begin{cases}
 2r,&w\leq-r/4,\\
 (12r-8w)/7,&-r/4<w<5r/8,\\
 r,&w\geq5r/8,
 \end{cases}
\]
with the sign of \(z_1\). The scale coordinate is respectively
\(8w\), \((-8r+24w)/7\), and \(8w/5\).
Thus every draw of the limit is specified by three arithmetic
branches and two independent normal variables.
Its location variance is approximately \(1.83270002\);
the annex gives an exact trigonometric expression.
The proposed variance estimate has limiting values
\(5/2,10/9,5/8\) on the corresponding regions.
\end{hbtheorem}

\begin{hbdefinition}{def:p15-huber-varying-cutoff-notation}{The inherited normal moment functions}
For a varying cutoff \(\ell>0\), write
\(\psi_\ell(x)=\max(-\ell,\min(x,\ell))\).
The fixed score above is \(\psi_k\).
Retain the normal moment functions of \cref{sec:p15-huber-table-one}:
\[
 A(\ell)=2\Phi(\ell)-1,\qquad
 B(\ell)=A(\ell)-2\ell\phi(\ell)+2\ell^2(1-\Phi(\ell)).
\]
The function \(A(\ell)\) is distinct from the bare interior probability
\(A=P(|U|<k)\).
\end{hbdefinition}

\begin{hbtheorem}{thm:p15-huber-joint-normal-contamination-counterexample}{proved=annex:p15.huber.joint.normal.contamination.counterexample}{A non-Gaussian location limit inside the original model}
The location marginal, not only the joint vector, can fail to be
Gaussian. Here is a counterexample with the original normal calibration.
Set \(k=1\), \(\beta=B(1)\), \(p=1/1000\), and \(a_0=1-2p\).
Let \(q\) be the unique root
\[
              B(q)/q^2=(\beta-2p)/a_0,
\]
and take
\[
                F=a_0\Phi+p\delta_{-q}+p\delta_q .
\]
Then \(1<q<2\), \(S_0=q\), and
\[
 A=a_0 A(q),\qquad
 C=a_0\{A(q)-2q\phi(q)\}/q^2 .
\]
Numerically \(q\simeq1.00244151\), \(A\simeq0.68250186\),
and \(C\simeq0.19856041\).
The resulting location limit is non-Gaussian.
The proof establishes this from its tail rate and second moment,
not from a simulated histogram.
\end{hbtheorem}

\begin{hbtheorem}{thm:p15-huber-joint-atomic-bound}{proved=annex:p15.huber.joint.atomic.bound}{The sharp source upper bound survives with its corrected interpretation}
For an identified symmetric normal-contamination law
\(F=(1-\varepsilon)\Phi+\varepsilon H_{\rm c}\), take
\(\beta=B(k)\), \(r=B(k)/k^2\), and \(0\leq\varepsilon<r\).
With \(q_{\max}\) defined in
\cref{thm:p15-huber-joint-variance-bound},
\[
                  S_0^2 EH_1^2\leq W(q_{\max},\varepsilon).
\]
This includes clipping atoms; it bounds the variance of the
possibly non-Gaussian limiting location law.
It is sharp over the identified laws, attained by symmetric
contaminants strictly outside the final cutoff.
If also \(\varepsilon<1-r\), every symmetric contaminant is
identified, so the same is the sharp bound over the entire
source contamination class. This gives the warranted interpretation
of Table IV without retaining an unjustified universal normal limit.
\end{hbtheorem}

\begin{hbtheorem}{thm:p15-huber-zero-atom-mixture-bridge}{proved=annex:p15.huber.zero.atom.mixture.bridge}{The normal-mixture comparison needs no contamination-fraction restriction}
For every symmetric law with \(F\{0\}=0\), and every
\(0<\beta<k^2\), the identification and consistency theorem applies.
If the law is non-atomic, \(p=0\), and the Gaussian special case
applies as well. In particular, for
\[
 F_\varepsilon(x)=(1-\varepsilon)\Phi(x)+\varepsilon\Phi(x/3),
                       \qquad 0\leq\varepsilon\leq1,
\]
the entire consistency and variance calculation in
\cref{thm:p15-huber-pair-median} is justified, with
\(\beta=B(1.5)\), without requiring
\(\varepsilon<1-B(1.5)/(1.5)^2\).
The earlier restriction was sufficient for arbitrary contaminants;
it was not a necessary property of these continuous mixtures.
\end{hbtheorem}

\section{All positive scale-equivariant estimators, including signs}
\label{sec:p15-huber-scale-signs}

\begin{hbdefinition}{def:p15-huber-scale-sign-class}{The exact comparison class and ordered risk}
Fix \(q\geq\sqrt2\) and the corresponding positive contamination
fraction and information \(I_Y\) of
\cref{thm:p15-huber-scale-model}.
Let \(\mathcal F_q\) consist of laws
\(F=(1-\varepsilon)\Phi+\varepsilon H\) with
\(H\{|X|>q\}=1\). No symmetry or independence of sign and magnitude
is imposed on \(H\).
A competing scale estimate \(S_n\) is any finite positive measurable
function on \((\mathbb R\setminus\{0\})^n\) satisfying
\[
                         S_n(cx)=cS_n(x)\qquad(c>0).
\]
It may use every sign and magnitude. This domain has probability one
under every law under consideration. There is no finite positive
scale-equivariant value at the all-zero vector, since equivariance
there would require \(S_n(0)=cS_n(0)\) for every \(c>0\).
Let \(\sigma F\) denote the law of \(\sigma X\), \(X\sim F\), and set
\[
 \mathcal R_q(S)=
 \sup_{F\in\mathcal F_q}\sup_{\sigma>0}
 \lim_{R\to\infty}\liminf_{n\to\infty}
 E_{\sigma F}\min\{n[\log S_n^2-\log\sigma^2]^2,R^2\}.
\]
The data in each expectation are i.i.d.
\end{hbdefinition}

\begin{hbtheorem}{thm:p15-huber-scale-sign-minimax}{proved=annex:p15.huber.scale.sign.minimax}{The finite sign-conditioning lower bound and full minimax equality}
There is a least-favorable \(F_*\in\mathcal F_q\) for which,
for every competing \(S_n\), every \(n\geq1\), and every \(R>0\),
\[
 E_{F_*}\min\{n(\log S_n^2)^2,R^2\}
                   \geq\frac1{I_Y+4\pi^2/R^2}.
\]
It has the explicit even density
\[
 f_*(x)=
 \begin{cases}
 (1-\varepsilon)\phi(x),&|x|\leq q,\\
 (1-\varepsilon)\phi(q)(|x|/q)^{-q^2},&|x|>q.
 \end{cases}
\]
The proof conditions on independent fair signs at this one
least-favorable law. Conditional on each sign pattern the original
estimator becomes a deterministic translation-equivariant estimator
of the log-squared observations; the comparison class is not changed.
No artificial randomization or Jensen step on truncated loss is used.
Consequently
\[
          \inf_{S\ {\rm positive\ scale\ equivariant}}\mathcal R_q(S)
                       =1/I_Y .
\]
The source's scale estimate, which uses the magnitudes, attains
this value. Scale equivariance makes the risk independent of
\(\sigma\), and the two limiting operations retain the displayed
order.
\end{hbtheorem}
